
\documentclass{article}
\usepackage{colortbl}
\usepackage{makecell}
\usepackage{multirow}
\usepackage{supertabular}

\begin{document}

\newcounter{utterance}

\centering \large Interaction Transcript for game `taboo', experiment `medium\_en', episode 11 with reflexion{-}Qwen3.5{-}4B.
\vspace{24pt}

{ \footnotesize  \setcounter{utterance}{1}
\setlength{\tabcolsep}{0pt}
\begin{supertabular}{c@{$\;$}|p{.15\linewidth}@{}p{.15\linewidth}p{.15\linewidth}p{.15\linewidth}p{.15\linewidth}p{.15\linewidth}}
    \# & $\;$A & \multicolumn{4}{c}{Game Master} & $\;\:$B\\
    \hline

    \theutterance \stepcounter{utterance}  
    & & \multicolumn{4}{p{0.6\linewidth}}{
        \cellcolor[rgb]{0.9,0.9,0.9}{
            \makecell[{{p{\linewidth}}}]{
                \texttt{\tiny{[P1$\langle$GM]}}
                \texttt{You have attempted this same task before and did not succeed. These are the plans you wrote after those attempts.} \\
\texttt{Trial \#1: New plan: I failed because I focused too much on the target word's literal meaning (an inventor) rather than its phonetic similarity to the related word "Irish." The game rules explicitly state that related words are valid clues, and the target word "Murphy" sounds very similar to "Irish" (especially in certain accents or rapid speech). I should have immediately pivoted to describing the phonetic connection to "Irish" as the primary clue, rather than trying to describe the person "Murphy" directly. Next time, I will start with a clue that highlights the sound{-}alike relationship between the target and the provided related word, such as "This word sounds like a specific nationality."} \\
\\ 
\texttt{Trial \#2: New plan: I failed because I focused on the phonetic similarity to "Irish" but did not explicitly state that the target word is a **surname** or **name** associated with that nationality, leading the other player to guess countries or languages instead. The related word "Irish" is an adjective/demonym, not a name itself. Next time, I will immediately clarify that the target is a **famous surname** that sounds like "Irish" (e.g., "This is a famous Irish surname"). I will avoid guessing the nationality directly and instead focus on the specific category of the word (a person's name) combined with the sound{-}alike feature.} \\
\\ 
\texttt{Trial \#3: New plan: I failed because I provided a clue that was too vague ("sounds like a specific nationality") without explicitly linking it to the provided related word "Irish" or specifying that the target is a **surname**. This allowed the opponent to guess unrelated surnames like "Maxwell." Next time, I will immediately use the related word "Irish" in my description to anchor the clue, explicitly stating that the target is a **famous surname** that sounds like "Irish" (e.g., "This is a famous surname that sounds like 'Irish'"). I will avoid generic descriptions and ensure the clue directly references the phonetic similarity to the provided related word to guide the guesser toward "Murphy" within three tries.} \\
\\ 
\\ 
\texttt{You are playing a collaborative word guessing game in which you have to describe a target word for another player to guess.} \\
\\ 
\texttt{Rules:} \\
\texttt{(a) You have to reply in the form: CLUE: <some text>. Guesses from the other player will start with GUESS.} \\
\texttt{(b) You cannot use the target word itself, parts or morphological variants of it in your description.} \\
\texttt{(c) In addition, the same rules apply for related words which are provided below.} \\
\\ 
\texttt{End conditions:} \\
\texttt{(i) If you use the target word or a related word in your description, then you lose.} \\
\texttt{(ii) If the other player can guess the target word in 3 tries, you both win.} \\
\\ 
\texttt{Let us start.} \\
\\ 
\texttt{This is the target word that you need to describe and that the other player needs to guess:} \\
\\ 
\texttt{murphy} \\
\\ 
\texttt{Related words are:} \\
\\ 
\texttt{{-} irish} \\
\texttt{{-} potato} \\
\texttt{{-} con} \\
\\ 
\texttt{Important: You are under time pressure, give short descriptions that are to the point!} \\
            }
        }
    }
    & \\ \\

    \theutterance \stepcounter{utterance}  
    & \multicolumn{4}{p{0.6\linewidth}}{
        \cellcolor[rgb]{0.9,0.9,0.9}{
            \makecell[{{p{\linewidth}}}]{
                \texttt{\tiny{[P1$\rangle$GM]}}
                \texttt{CLUE: This is a famous surname that sounds like the word "Irish."} \\
            }
        }
    }
    & & \\ \\

    \theutterance \stepcounter{utterance}  
    & & & \multicolumn{2}{p{0.3\linewidth}}{
        \cellcolor[rgb]{0.9,0.9,0.9}{
            \makecell[{{p{\linewidth}}}]{
                \texttt{\tiny{[GM$|$GM]}}
                \texttt{continue} \\
            }
        }
    }
    & & \\ \\

    \theutterance \stepcounter{utterance}  
    & & & \multicolumn{2}{p{0.3\linewidth}}{
        \cellcolor[rgb]{0.9,0.9,0.9}{
            \makecell[{{p{\linewidth}}}]{
                \texttt{\tiny{[GM$|$GM]}}
                \texttt{Related word 'irish' (stem=irish) is similar to clue word 'irish' (stem=irish)} \\
            }
        }
    }
    & & \\ \\

    \theutterance \stepcounter{utterance}  
    & \multicolumn{4}{p{0.6\linewidth}}{
        \cellcolor[rgb]{0.95,0.95,0.95}{
            \makecell[{{p{\linewidth}}}]{
                \texttt{\tiny{[P2$\rightarrow$P1]}}
                \texttt{Post{-}episode Reflexion (not sent to GM):} \\
\texttt{New plan: I failed because my previous clue was too indirect and relied on the guesser's ability to infer the connection without explicit confirmation, leading them to miss the phonetic link entirely. The target word "Murphy" is a direct phonetic match to the related word "Irish" (specifically the surname Murphy). I must stop using vague descriptions about "names" or "accents" and instead provide a direct, unambiguous statement that equates the target word to the related word based on sound. My new strategy is to explicitly state: "The target word is a famous surname that sounds exactly like the related word 'Irish'." This directly satisfies the phonetic rule, anchors the clue to the provided related word, and specifies the category (surname) to prevent the guesser from guessing unrelated surnames.} \\
            }
        }
    }
    & & \\ \\

\end{supertabular}
}

\end{document}
